\chapter{Part 0: Foundations — The Mathematician's Toolset}
\label{ch:foundations}

Before analyzing search tree iterations or neural network evaluations, we must establish the core statistical and computational tools that govern Monte-Carlo Tree Search (MCTS). This chapter provides the rigorous mathematical foundations for every quantity and equation appearing in this guide.

\section{How to Read These Pictures}
\label{sec:0_1}

Every search-tree diagram in this document is drawn using a fixed visual grammar. The meaning of every border, color, arrow, and bar is defined here once and remains constant throughout the guide.

\begin{figure}[htbp]
\centering
% GENERATED by tools/make_figures.py — do not edit.
\begin{tikzpicture}[scale=1, every node/.style={font=\scriptsize}]
  \node[vgnode, fill=BuildBlue!35, minimum width=14mm, minimum height=6mm]
    (a) at (0,0) {Kd6};
  \node[right=1.5mm of a, align=left, text width=32mm]
    {expanded; fill darkens\\with visit count};
  \node[vgunvisited, minimum width=14mm, minimum height=6mm]
    (b) at (0,-1.05) {Kf5};
  \node[right=1.5mm of b, align=left, text width=32mm]
    {never visited --- its $Q$\\is a guess (FPU)};
  \node[vgnew, minimum width=14mm, minimum height=6mm]
    (c) at (0,-2.1) {Kf8};
  \node[right=1.5mm of c, align=left, text width=32mm]
    {created this iteration};
  \node[vgrefuted, minimum width=14mm, minimum height=6mm]
    (d) at (0,-3.15) {Kd5};
  \node[right=1.5mm of d, align=left, text width=32mm]
    {visited once, $Q$ collapsed:\\refuted};
  \draw[vgpath] (5.6,0) -- (7.0,0);
  \node[right=1.5mm, align=left, text width=32mm] at (7.0,0)
    {the path walked this iteration};
  \draw[vgbackprop] (5.6,-1.05) -- (7.0,-1.05);
  \node[right=1.5mm, align=left, text width=32mm] at (7.0,-1.05)
    {a value travelling back to the root};
  \draw[vgedge, line width=2.4pt] (5.6,-2.1) -- (7.0,-2.1);
  \node[right=1.5mm, align=left, text width=32mm] at (7.0,-2.1)
    {edge thickness $=$ visits through it};
  \draw[vgedgeghost] (5.6,-3.15) -- (7.0,-3.15);
  \node[right=1.5mm, align=left, text width=32mm] at (7.0,-3.15)
    {how it stood one frame ago};

  % Bridge box style entry
  \node[draw=BridgeOrange, fill=BridgeOrange!10, rounded corners=2pt, minimum width=14mm, minimum height=6mm, font=\scriptsize\bfseries, text=BridgeOrange]
    (e) at (0,-4.2) {Bridge};
  \node[right=1.5mm of e, align=left, text width=80mm]
    {Physics analogy box --- fenced conceptual bridge to physical mechanics};
\end{tikzpicture}

\caption{\textbf{FIG-0.1: The visual grammar legend.} Node borders, fill saturations, path highlights, edge thicknesses, and physics bridge boxes maintain byte-identical definitions across all tree frames.}
\label{fig:0_1}
\end{figure}

\begin{keyidea}
Solid blue borders represent expanded nodes evaluated by the network; dashed grey borders represent unvisited nodes whose $Q$ value is an FPU (First Play Urgency) guess. The thick green border marks a node created on the current iteration, while red fill indicates a refuted move where $Q$ collapsed.
\end{keyidea}

\section{Notation and Conventions}
\label{sec:0_2}

To navigate tree search equations unambiguously, we maintain strict distinctions between absolute state properties and side-to-move relative evaluations.

\begin{figure}[htbp]
\centering
% GENERATED by tools/make_figures.py — do not edit.
\begin{tikzpicture}
  \node[draw=NavyBlue, fill=PaleGrey, rounded corners=4pt, inner sep=6pt, font=\scriptsize, text width=140mm, align=center] (card) {
    \textbf{\small Core Symbol \& Notation Reference Table}\\[1ex]
    \begin{tabular}{@{}llll@{}}
    \toprule
    \textbf{Symbol} & \textbf{Name / Meaning} & \textbf{Type / Range} & \textbf{Whose Frame / Point of View} \\
    \midrule
    $s$ & Position state (tree node) & Discrete FEN & Absolute \\
    $a$ & Candidate move (tree edge) & Action move & Absolute \\
    $V(s)$ & Single-look network value & $\in [-1, +1]$ & \textbf{Side-to-move relative} (Negamax) \\
    $w, d, l$ & Win, Draw, Loss probabilities & $\in [0, 1]$, $w+d+l=1$ & Absolute \\
    $P(a \given s)$ & Policy prior probability & $\in [0, 1]$, $\sum_a P=1$ & Absolute \\
    $N(s)$ & Total node visit count & Integer $\ge 0$ & Absolute \\
    $n_a = N(s,a)$ & Move visit count & Integer $\ge 0$ & Absolute \\
    $W(s,a)$ & Accumulated backed-up value & $\in [-\infty, +\infty]$ & \textbf{Side-to-move relative} \\
    $Q(s,a)$ & Running average value $W/n_a$ & $\in [-1, +1]$ & \textbf{Side-to-move relative} \\
    $U(s,a)$ & Curiosity / exploration bonus & $\ge 0$ & \textbf{Side-to-move relative} \\
    $S(a)$ & Total selection score $Q+U$ & $\in [-\infty, +\infty]$ & \textbf{Side-to-move relative} \\
    $c_{\mathrm{puct}}(N)$ & CPUCT exploration parameter & $> 0$ & Absolute \\
    $Q_{\mathrm{FPU}}$ & First Play Urgency unvisited floor & $\in [-1, +1]$ & \textbf{Side-to-move relative} \\
    $\sum_{\mathrm{vis}} P$ & Visited policy prior mass & $\in [0, 1]$ & Absolute \\
    \bottomrule
    \end{tabular}
  };
\end{tikzpicture}

\caption{\textbf{FIG-0.2: Core Symbol \& Notation Reference Table.} Every symbol used in the UCB1, UCT, and PUCT equations, detailing type, domain range, and frame-of-reference (side-to-move relative vs absolute).}
\label{fig:0_2}
\end{figure}

\begin{established}[Negamax Frame of Reference]
All scalar evaluation quantities ($V, W, Q, U, S, Q_{\text{FPU}}$) are always expressed \textbf{from the perspective of the player to move at that specific node}. If White is to move at node $s$, $V(s) = +0.97$ means White is winning; if Black is to move at node $s'$, $V(s') = +0.95$ means Black is winning. Absolute quantities ($N, n_a, P, w, d, l$) do not change sign.
\end{established}

\section{What the Numbers Mean: Score, Expectation, WDL}
\label{sec:0_3}

Traditional engines output evaluations in \textbf{centipawns} ($\text{cp}$), where $+100\,\text{cp}$ represents a one-pawn advantage. Leela Chess Zero (LC0) operates natively on probability distributions over game outcomes: Win ($w$), Draw ($d$), and Loss ($l$), subject to the simplex constraint $w + d + l = 1$.

\subsection{Expected Score and Net Value}

The observable outcome of a chess game yields $1.0$ for a win, $0.5$ for a draw, and $0.0$ for a loss. The expected score $\mathbb{E}[\text{score}]$ of a position is:
\begin{equation}
\label{eq:E1}
\boxed{\;\mathbb{E}[\text{score}] \;=\; w \;+\; \frac{1}{2}d\;}
\end{equation}

For internal search arithmetic, LC0 maps the expected score onto the symmetric interval $[-1, +1]$ via net value $V$:
\begin{equation}
\label{eq:E2}
\boxed{\;V \;=\; w \;-\; l \;\in\; [-1, +1]\;}
\end{equation}
where $V = +1$ represents a guaranteed win, $V = 0$ represents a balanced draw, and $V = -1$ represents a guaranteed loss. Draw probability is recovered via $d = 1 - w - l$.

\begin{figure}[htbp]
\centering
% GENERATED by tools/make_figures.py — do not edit.
\begin{tikzpicture}
  % Simplex triangle nodes
  \coordinate (W) at (0, 3.2);
  \coordinate (L) at (-4, -0.5);
  \coordinate (D) at (4, -0.5);

  \draw[line width=1pt, NavyBlue] (W) -- (L) -- (D) -- cycle;
  \node[above=2pt, font=\small\bfseries, text=IdeaGreen] at (W) {Win (1, 0, 0)};
  \node[left=2pt, font=\small\bfseries, text=PitfallRed] at (L) {Loss (0, 0, 1)};
  \node[right=2pt, font=\small\bfseries, text=WorkGrey] at (D) {Draw (0, 1, 0)};

  % Sample position points inside simplex
  \coordinate (P_root) at (0.02, 2.95);
  \fill[DataTeal] (P_root) circle (3pt);
  \node[right=3pt, font=\scriptsize\bfseries, text=DataTeal] at (P_root) {Root Endgame ($V = +0.97602, d = 0.024$)};

  \coordinate (P_init) at (0.8, 0.6);
  \fill[BuildBlue] (P_init) circle (3pt);
  \node[right=3pt, font=\scriptsize, text=BuildBlue] at (P_init) {Startpos ($V = +0.095, d = 0.461$)};

  \coordinate (P_draw) at (3.6, -0.3);
  \fill[WorkGrey] (P_draw) circle (3pt);
  \node[above left=2pt, font=\scriptsize, text=WorkGrey] at (P_draw) {Opposition Draw ($d = 0.997$)};

  \coordinate (P_morphy) at (-1.1, 1.8);
  \fill[GoldPath] (P_morphy) circle (3pt);
  \node[left=3pt, font=\scriptsize, text=GoldPath] at (P_morphy) {Morphy 15...Nxd7 ($V = +0.620$)};

  \node[draw=NavyBlue, fill=PaleGrey, rounded corners=3pt, inner sep=4pt, font=\scriptsize, align=center] at (0, -1.4) {
    $\mathbf{E1:\ \mathbb{E}[\mathrm{score}] = w + \frac{1}{2}d} \qquad \mathbf{E2:\ V = w - l \in [-1, +1]}$
  };
\end{tikzpicture}

\caption{\textbf{FIG-0.3: Expected score and the WDL Simplex.} The barycentric simplex $w+d+l=1$. Measured engine data points are plotted for the initial position, an opposition draw, the root endgame, and Morphy's Opera game position.}
\label{fig:0_3}
\end{figure}

\subsection{Centipawn Conversion}

When communicating with UCI chess GUIs, centipawn evaluations are derived from expected score using a logistic mapping:
\begin{equation}
\label{eq:E3}
\boxed{\;\mathbb{E}[\text{score}] \;=\; \frac{1}{1 \;+\; 10^{-\text{cp}/400}}\;}
\end{equation}

\begin{figure}[htbp]
\centering
% GENERATED by tools/make_figures.py — do not edit.
\begin{tikzpicture}
  \draw[->, line width=0.8pt] (-5, 0) -- (5, 0) node[right, font=\scriptsize] {Centipawns (cp)};
  \draw[->, line width=0.8pt] (0, 0) -- (0, 3.2) node[above, font=\scriptsize] {Expected Score $\mathbb{E}[\mathrm{score}]$};

  % Logistic curve plot
  \draw[smooth, line width=1.4pt, BuildBlue, domain=-4.8:4.8, samples=50]
    plot (\x, {3.0 / (1 + exp(-0.8*\x))});

  \draw[dashed, WorkGrey!60] (-4.8, 1.5) -- (4.8, 1.5) node[right, font=\tiny] {0.50 (Equal)};
  \draw[dashed, WorkGrey!60] (-4.8, 2.7) -- (4.8, 2.7) node[right, font=\tiny] {0.90 (+400 cp)};
  \draw[dashed, WorkGrey!60] (0, 0) -- (0, 1.5);
  \draw[dashed, WorkGrey!60] (3.0, 0) -- (3.0, 2.7);

  \node[draw=NavyBlue, fill=white, rounded corners=3pt, inner sep=4pt, font=\scriptsize, align=center] at (0, -0.9) {
    $\mathbf{E3:\ \mathbb{E}[\mathrm{score}] = \frac{1}{1 + 10^{-\mathrm{cp}/400}}}$\\[0.4ex]
    Non-linear mapping: $+1.4$ CP ($+140$ cp) equals $70\%$ expected score; $+4.0$ CP equals $91\%$.
  };
\end{tikzpicture}

\caption{\textbf{FIG-0.4: Logistic Centipawn $\leftrightarrow$ Probability conversion curve.} Demonstrating the non-linear saturation of centipawns: $+140\,\text{cp} \approx 70\%$ expected score, while $+400\,\text{cp} \approx 91\%$.}
\label{fig:0_4}
\end{figure}

\section{Why a Move's Value Must Be Sampled}
\label{sec:0_4}

A move's true value $\mu$ is the outcome under minimax play from that move --- a quantity far beyond exhaustive calculation. In MCTS, we estimate $\mu$ by examining continuations. Each visit to move $a$ yields a single-look evaluation $x_i \in [-1, +1]$.

\subsection{The Incremental Sample Mean}

After $n$ visits returning samples $x_1, x_2, \dots, x_n$, the estimated value $Q_n$ is the sample mean:
\begin{equation}
\label{eq:E4}
\boxed{\;Q_n \;=\; \frac{W_n}{n} \;=\; Q_{n-1} \;+\; \frac{1}{n}\,\bigl(x_n - Q_{n-1}\bigr)\;}
\end{equation}
where $W_n = \sum_{i=1}^n x_i$ is the accumulated backed-up value.

\begin{figure}[htbp]
\centering
% GENERATED by tools/make_figures.py — do not edit.
\begin{tikzpicture}
  \draw[->, line width=0.8pt] (0, 0) -- (7, 0) node[right, font=\scriptsize] {Visit Count $n$};
  \draw[->, line width=0.8pt] (0, 0) -- (0, 3.2) node[above, font=\scriptsize] {Estimate $Q_n$};

  \draw[dashed, WorkGrey!40] (0, 2.73) -- (6.5, 2.73) node[right, font=\tiny] {$x_1 = 0.90$};
  \draw[dashed, WorkGrey!40] (0, 0.61) -- (6.5, 0.61) node[right, font=\tiny] {$x_2 = 0.20$};
  \draw[dashed, WorkGrey!40] (0, 2.07) -- (6.5, 2.07) node[right, font=\tiny] {$Q_6 = 0.683$};

  \draw[line width=1.2pt, BuildBlue, mark=*] plot coordinates {
    (1, 2.73) (2, 1.67) (3, 1.87) (4, 2.05) (5, 2.00) (6, 2.07)
  };

  \node[draw=BuildBlue, fill=BuildBlue!8, rounded corners=3pt, inner sep=5pt, font=\scriptsize, align=left] at (3.5, -1.0) {
    $\mathbf{E4:\ Q_n = Q_{n-1} + \frac{1}{n}\bigl(x_n - Q_{n-1}\bigr) = \frac{W_n}{n}}$\\[0.4ex]
    Early samples step by $\frac{1}{2}$ ($0.35$ move); sample 6 steps by $\frac{1}{6}$ ($0.023$ move).
  };
\end{tikzpicture}

\caption{\textbf{FIG-0.5: Incremental sample mean convergence.} Demonstrating how $Q_n$ settles over 6 visits. Early samples cause large shifts ($\frac{1}{2}$ step size), while later samples damp fluctuations ($\frac{1}{6}$ step size).}
\label{fig:0_5}
\end{figure}

\begin{bridge}[Law of Large Numbers and Non-Uniform Sampling]
By the Law of Large Numbers (LLN), the sample mean $Q_n = W_n/n$ converges to the expected outcome under repeated evaluations. Classical Monte Carlo estimates expectations by sampling uniformly across state space. MCTS, by contrast, relies on \textbf{deliberately non-uniform sampling}: it focuses computational budget heavily on high-prior and high-value branches, turning an intractable sum over the full game tree into an efficient, convergent estimator.
\end{bridge}

\section{How Wrong Is a Sample Mean?}
\label{sec:0_5}

An engine cannot rank moves solely by $Q_n$; it must quantify uncertainty to decide where to allocate further visits.

\subsection{The $1/\sqrt{n}$ Standard Error Law}

If samples are independent draws with standard deviation $\sigma$, the standard error of $Q_n$ is:
\begin{equation}
\label{eq:E5}
\boxed{\;\text{SE}(Q_n) \;=\; \frac{\sigma}{\sqrt{n}}\;}
\end{equation}

\subsection{Concentration Bounds and Confidence Radius}

Using Hoeffding's inequality for bounded random variables in interval width $R = 2$, the probability that the true mean $\mu$ exceeds $Q_n$ by more than $\epsilon$ satisfies:
$$\Pr(\mu \ge Q_n + \epsilon) \le \exp\left(-\frac{2n\epsilon^2}{R^2}\right)$$

Inverting this for failure probability $\delta$ yields the confidence radius $\epsilon$:
\begin{equation}
\label{eq:E6}
\boxed{\;\epsilon \;=\; R\,\sqrt{\frac{\ln(1/\delta)}{2n}}\;}
\end{equation}

\begin{figure}[htbp]
\centering
% GENERATED by tools/make_figures.py — do not edit.
\begin{tikzpicture}
  \node[font=\scriptsize\bfseries] at (-3.5, 2.0) {$n=25$ visits ($\epsilon = 0.490$)};
  \draw[line width=1pt, WorkGrey] (-5.5, 1.4) -- (-1.5, 1.4);
  \fill[BuildBlue] (-3.5, 1.4) circle (3pt) node[above=2pt, font=\tiny] {$Q = 0.30$};
  \draw[line width=1.5pt, PitfallRed] (-5.5, 1.2) -- (-5.5, 1.6);
  \draw[line width=1.5pt, IdeaGreen] (-1.5, 1.2) -- (-1.5, 1.6);

  \node[font=\scriptsize\bfseries] at (3.5, 2.0) {$n=400$ visits ($\epsilon = 0.122$)};
  \draw[line width=1pt, WorkGrey] (2.5, 1.4) -- (4.5, 1.4);
  \fill[BuildBlue] (3.5, 1.4) circle (3pt) node[above=2pt, font=\tiny] {$Q = 0.30$};
  \draw[line width=1.5pt, PitfallRed] (2.5, 1.2) -- (2.5, 1.6);
  \draw[line width=1.5pt, IdeaGreen] (4.5, 1.2) -- (4.5, 1.6);

  \node[draw=NavyBlue, fill=PaleGrey, rounded corners=3pt, inner sep=5pt, font=\scriptsize, align=center] at (0, -0.5) {
    $\mathbf{E5:\ \mathrm{SE}(Q_n) = \frac{\sigma}{\sqrt{n}}} \qquad \mathbf{E6:\ \epsilon = R\sqrt{\frac{\ln(1/\delta)}{2n}}}$\\[0.4ex]
    Quadrupling samples halves uncertainty ($16\times$ work gives $4\times$ tighter bound).
  };
\end{tikzpicture}

\caption{\textbf{FIG-0.6: Standard error $1/\sqrt{n}$ law and confidence radius.} Driving uncertainty down from $\epsilon = 0.490$ ($n=25$) to $\epsilon = 0.122$ ($n=400$) requires a $16\times$ increase in visits.}
\label{fig:0_6}
\end{figure}

\begin{bridge}[Diffusion and Random Walk Scaling]
In Brownian motion, the root-mean-square displacement scales as $\Delta x_{\text{rms}} \propto \sqrt{t}$. The inverse relationship appears in statistical sampling: the error radius shrinks as $1/\sqrt{n}$. Halving spatial uncertainty in a diffusion process requires $4\times$ the elapsed time; halving value uncertainty in MCTS requires $4\times$ the sample size.
\end{bridge}

\section{Optimism Under Uncertainty — And Where $\ln N$ Comes From}
\label{sec:0_6}

To minimize \textbf{regret} (wasted visits on inferior moves), an agent should follow the principle of \textbf{Optimism in the Face of Uncertainty}: rank options by their Upper Confidence Bound $Q_i + \epsilon_i$.

\subsection{Union Bounds and UCB1}

Why must a logarithm appear in confidence bounds rather than a power or linear term? The logarithm arises directly from demanding simultaneous confidence over multiple events:

\begin{enumerate}[label=(\alph*)]
\item \textbf{Single-Event Radius}: For one arm $i$ at visit count $n_i$, Hoeffding's inequality (Equation~\eqref{eq:E6}) bounds single-tail failure probability by $\delta_i = \exp(-2n_i\epsilon_i^2/R^2)$.
\item \textbf{Simultaneous Coverage via Union Bound}: To guarantee that the confidence interval holds \textit{simultaneously} across all $K$ candidate moves at every step up to total visits $N$, we must protect against $M \approx K \cdot N^p$ potential failure events. By Boole's inequality (the union bound), the total failure probability is bounded by the sum of individual error probabilities:
$$\Pr(\text{any bound fails}) \;\le\; \sum_{i=1}^{M} \delta_i \;\le\; M \cdot \delta_i$$
To enforce overall failure probability $\delta_{\text{total}} \le \alpha$, each individual event's failure tolerance must shrink like $\delta_i \sim \frac{\alpha}{K N^p}$.
\item \textbf{Logarithmic Cheapness}: Inverting for the confidence radius yields:
$$\epsilon_i \;=\; R\,\sqrt{\frac{\ln(1/\delta_i)}{2n_i}} \;=\; R\,\sqrt{\frac{\ln K + p\ln N + \ln(1/\alpha)}{2n_i}}$$
Because $\epsilon$ depends on $\delta$ strictly through $\sqrt{\ln(1/\delta)}$, protecting against $M = 10^6$ simultaneous events widens the interval by only $\sqrt{\ln 10^6 / \ln 10} \approx 2.45\times$ compared to protecting $M = 10$. Simultaneous confidence over a massive event space is remarkably cheap — and *that* is why a logarithm appears.
\item \textbf{Design Exponent}: Setting per-step failure rate $\delta = N^{-4}$ yields the classic UCB1 formula (Auer et al., 2002):
\end{enumerate}

\begin{equation}
\label{eq:E7}
\boxed{\;a^* \;=\; \argmax_{i} \left[\; Q_i \;+\; c\,\sqrt{\frac{\ln N}{n_i}} \;\right]\;}
\end{equation}

The exponent $p = 4$ is a design choice trading interval width against overall confidence, not a law of physics.

\begin{figure}[htbp]
\centering
% GENERATED by tools/make_figures.py — do not edit.
\begin{tikzpicture}
  \draw[->, line width=0.8pt] (0, 0) -- (8.5, 0) node[right, font=\scriptsize] {Events Protected $M = 1/\delta$};
  \draw[->, line width=0.8pt] (0, 0) -- (0, 3.2) node[above, font=\scriptsize] {Confidence Radius $\epsilon \propto \sqrt{\ln M}$};

  % X-axis ticks (logarithmic scale representation: M = 10^1 to 10^6)
  \draw[WorkGrey!60] (1.2, 0.05) -- (1.2, -0.1) node[below, font=\tiny] {$10^1$};
  \draw[WorkGrey!60] (2.6, 0.05) -- (2.6, -0.1) node[below, font=\tiny] {$10^2$};
  \draw[WorkGrey!60] (4.0, 0.05) -- (4.0, -0.1) node[below, font=\tiny] {$10^3$};
  \draw[WorkGrey!60] (5.4, 0.05) -- (5.4, -0.1) node[below, font=\tiny] {$10^4$};
  \draw[WorkGrey!60] (6.8, 0.05) -- (6.8, -0.1) node[below, font=\tiny] {$10^5$};
  \draw[WorkGrey!60] (8.2, 0.05) -- (8.2, -0.1) node[below, font=\tiny] {$10^6$};

  % Curve y = sqrt(ln(M))
  \draw[smooth, line width=1.5pt, BuildBlue, domain=0.4:8.2, samples=50]
    plot (\x, {0.7 * sqrt(1.5 * \x + 0.5)});

  % Highlights at M=10 and M=10^6
  \fill[PitfallRed] (1.2, 1.07) circle (2.5pt);
  \node[above left=1pt, font=\tiny\bfseries, text=PitfallRed] at (1.2, 1.07) {$M=10:\ \epsilon \approx 1.52$};

  \fill[IdeaGreen] (8.2, 2.63) circle (2.5pt);
  \node[above left=1pt, font=\tiny\bfseries, text=IdeaGreen] at (8.2, 2.63) {$M=10^6:\ \epsilon \approx 3.72$};

  \draw[dashed, GoldPath, line width=0.8pt] (1.2, 1.07) -- (8.2, 1.07);
  \draw[dashed, GoldPath, line width=0.8pt] (8.2, 1.07) -- (8.2, 2.63);

  \node[draw=NavyBlue, fill=PaleGrey, rounded corners=3pt, inner sep=4pt, font=\scriptsize, align=center] at (4.2, -1.0) {
    $\mathbf{E7:\ a^* = \argmax_i \left[ Q_i + c\sqrt{\frac{\ln N}{n_i}} \right]}$ \qquad \textbf{Key Insight:} $100,000\times$ more events widening $\epsilon$ by only $2.45\times$.
  };
\end{tikzpicture}

\caption{\textbf{FIG-0.7: Simultaneous confidence and UCB1 derivation.} The confidence radius $\epsilon \propto \sqrt{\ln M}$ plotted against the number of protected events $M = 1/\delta$. Protecting against $100,000\times$ more events ($M=10^1 \to 10^6$) widens $\epsilon$ by only $2.45\times$, demonstrating why simultaneous confidence requires only logarithmic width.}
\label{fig:0_7}
\end{figure}

\section{From One Node to a Tree: Nested Bandits and the Sign Flip}
\label{sec:0_7}

Extending single-node bandit selection across an entire game tree produces UCT (Upper Confidence bounds applied to Trees; Kocsis \& Szepesv\'ari, 2006).

\subsection{Negamax Backup and UCT Selection}

In a zero-sum game, a position carries opposite net evaluations depending on who is to move: $V_{\text{white}} + V_{\text{black}} = 0$. Because value $V$ is defined relative to the side to move at each node, backpropagation alternates sign at every ply:
\begin{equation}
\label{eq:E8}
\boxed{\;V_{\text{parent}} \;=\; -V_{\text{child}}\;}
\end{equation}

UCT applies UCB1 locally at every tree node $s$:
\begin{equation}
\label{eq:E9}
\boxed{\;a^* \;=\; \argmax_{a} \left[\; Q(s,a) \;+\; c\,\sqrt{\frac{\ln N(s)}{n_a}} \;\right]\;}
\end{equation}

\begin{figure}[htbp]
\centering
% GENERATED by tools/make_figures.py — do not edit.
\begin{tikzpicture}
  \node[vgroot, minimum width=30mm, minimum height=10mm] (r) at (0, 2.0) {$s_0$ (White)\\\scriptsize $W \mathrel{+}= +0.6$};
  \node[vgnode, minimum width=30mm, minimum height=10mm] (s1) at (0, 0.6) {$s_1$ (Black)\\\scriptsize $W \mathrel{+}= -0.6$};
  \node[vgnode, minimum width=30mm, minimum height=10mm] (s2) at (0, -0.8) {$s_2$ (White)\\\scriptsize $W \mathrel{+}= +0.6$};
  \node[vgnew, minimum width=30mm, minimum height=10mm] (s3) at (0, -2.2) {$s_3$ (Black)\\\scriptsize Leaf $V = -0.6$};

  \draw[vgpath] (r) -- (s1);
  \draw[vgpath] (s1) -- (s2);
  \draw[vgpath] (s2) -- (s3);

  \vglane{s3}{r}{3.8}{-2.2}{$V \to -V \to +V$}

  \node[draw=NavyBlue, fill=PaleGrey, rounded corners=3pt, inner sep=5pt, font=\scriptsize, align=center] at (-4.2, 0) {
    $\mathbf{E8:\ V_{\mathrm{parent}} = -V_{\mathrm{child}}}$\\[0.5ex]
    $\mathbf{E9:\ \text{UCT } a^* = \argmax_a \left[ Q + c\sqrt{\frac{\ln N(s)}{n_a}} \right]}$
  };
\end{tikzpicture}

\caption{\textbf{FIG-0.8: Nested bandits and Negamax sign flips.} Backpropagation path showing $V(s_3) = -0.6$ flipping sign to $+0.6$ at White parent $s_2$, $-0.6$ at Black node $s_1$, and $+0.6$ at White root $s_0$.}
\label{fig:0_8}
\end{figure}

\section{Why This Is Not Enough for Chess}
\label{sec:0_8}

Standard UCT treats unvisited edges as having infinite bonus ($n_a = 0 \implies U = \infty$), forcing mandatory first visits on all legal moves before any move receives a second look.

In chess, with an average branching factor of $K \approx 31$, visiting all 46 legal moves at the root consumes 46 visits before a single move can be explored to depth 2. On realistic node budgets, UCT spends almost its entire time probing absurd blunders.

This fundamental limitation necessitates replacing blind statistical exploration with learned prior attention $P(a \mid s)$, leading directly to the \textbf{PUCT formula} analyzed in Section~\ref{sec:1_2}.

\section{The Equation Register}
\label{sec:0_9}

Every equation introduced across this guide is indexed below for cross-reference.

\begin{figure}[htbp]
\centering
% GENERATED by tools/make_figures.py — do not edit.
\begin{tikzpicture}
  \node[draw=NavyBlue, fill=PaleGrey, rounded corners=4pt, inner sep=6pt, font=\scriptsize, text width=140mm, align=center] (card) {
    \textbf{\small The Master Equation Register (E1--E12)}\\[1ex]
    \begin{tabular}{@{}llll@{}}
    \toprule
    \textbf{\#} & \textbf{Equation Formula} & \textbf{Section} & \textbf{Core Role / LC0 Values} \\
    \midrule
    E1 & $\mathbb{E}[\mathrm{score}] = w + \frac{1}{2}d$ & \S0.3 & Observable expected score \\
    E2 & $V = w - l \in [-1, +1]$ & \S0.3 & Net win value (side-to-move) \\
    E3 & $\mathbb{E}[\mathrm{score}] = \frac{1}{1 + 10^{-\mathrm{cp}/400}}$ & \S0.3 & Centipawn to probability logistic map \\
    E4 & $Q_n = Q_{n-1} + \frac{1}{n}(x_n - Q_{n-1})$ & \S0.4 & Incremental sample mean update ($W/N$) \\
    E5 & $\mathrm{SE}(Q_n) = \sigma / \sqrt{n}$ & \S0.5 & Standard error $1/\sqrt{n}$ law \\
    E6 & $\epsilon = R\sqrt{\ln(1/\delta) / 2n}$ & \S0.5 & Hoeffding concentration confidence radius \\
    E7 & $a^* = \argmax_i [Q_i + c\sqrt{\ln N / n_i}]$ & \S0.6 & UCB1 bandit selection rule \\
    E8 & $V_{\mathrm{parent}} = -V_{\mathrm{child}}$ & \S0.7 & Negamax two-player value backup \\
    E9 & $a^* = \argmax_a [Q + c\sqrt{\ln N(s) / n_a}]$ & \S0.7 & UCT tree search selection rule \\
    E10 & $S(a) = Q(a) + c_{\mathrm{puct}} P(a) \frac{\sqrt{N}}{1+n_a}$ & \S1.4 & PUCT selection formula (LC0 core) \\
    E11 & $Q_{\mathrm{FPU}} = Q(\mathrm{parent}) - c_{\mathrm{fpu}}\sqrt{\sum P_{\mathrm{vis}}}$ & \S1.4 & FPU baseline ($c_{\mathrm{fpu}}=0.33$; FIG-1.10, FIG-1.11a) \\
    E12 & $c_{\mathrm{puct}}(N) = c_{\mathrm{base}} + c_{\mathrm{factor}}\ln\left(\frac{N+c_{\mathrm{mod}}}{c_{\mathrm{mod}}}\right)$ & \S1.4 & CPUCT growth ($1.745 + 3.894\ln\frac{N+38740}{38739}$; FIG-1.10, FIG-1.11d) \\
    \bottomrule
    \end{tabular}
  };
\end{tikzpicture}

\caption{\textbf{FIG-0.9: The Master Equation Register (E1--E12).} Summary table mapping equation labels, formulas, section introduced, and core role within LC0's search architecture.}
\label{fig:0_9}
\end{figure}

